\RequirePackage[l2tabu, orthodox]{nag}
\documentclass[english, french]{beamer}
%INSTALL

%avoids a warning
\usepackage[log-declarations=false]{xparse}
\usepackage{fontspec} %font selecting commands
\usepackage{xunicode}
%warn about missing characters
\tracinglostchars=2

%REDAC
\usepackage{booktabs}
\usepackage{calc}

\usepackage{mathtools} %load this before babel!
	\mathtoolsset{showonlyrefs,showmanualtags}

\usepackage{babel}
%suppresses the warning about frenchb not modifying the captions (“—” to “:” in “Figure 1 – Legend”).
	\frenchbsetup{AutoSpacePunctuation=false,SuppressWarning=false}

\usepackage[super]{nth}
\usepackage{listings} %typeset source code listings
	\lstset{language=XML,tabsize=2,literate={"}{{\tt"}}1,captionpos=b}
\usepackage[nolist,smaller,printonlyused]{acronym}%,smaller option produces warnings from relsize in some cases, it seems.
\usepackage[nodayofweek]{datetime}%must be loaded after the babel package
\usepackage{xspace}
\usepackage{hyperref}% option pdfusetitle must be introduced here, not in hypersetup.
%breaklinks makes links on multiple lines into different PDF links to the same target.
%colorlinks (false): Colors the text of links and anchors. The colors chosen depend on the the type of link. In spite of colored boxes, the colored text remains when printing.
%linkcolor=black: this leaves other links in colors, e.g. refs in green, don't print well.
%pdfborder (0 0 1, set to 0 0 0 if colorlinks): width of PDF link border
%hidelinks
\hypersetup{breaklinks,bookmarksopen,colorlinks=true,urlcolor=blue,linkcolor=,hyperfigures=true}
% hyperref doc says: Package bookmark replaces hyperref’s bookmark organization by a new algorithm (...) Therefore I recommend using this package.
\usepackage{bookmark}

% center floats by default, but do not use with float
% \usepackage{floatrow}
% \makeatletter
% \g@addto@macro\@floatboxreset\centering
% \makeatother
\usepackage{ragged2e} %new com­mands \Cen­ter­ing, \RaggedLeft, and \RaggedRight and new en­vi­ron­ments Cen­ter, FlushLeft, and FlushRight, which set ragged text and are eas­ily con­fig­urable to al­low hy­phen­ation (the cor­re­spond­ing com­mands in LaTeX, all of whose names are lower-case, pre­vent hy­phen­ation al­to­gether). 
\usepackage{siunitx} %[expproduct=tighttimes, decimalsymbol=comma]
\sisetup{detect-all}% to detect e.g. when in math mode (use a math font)
\usepackage{braket} %for \Set
\usepackage{natbib}

\usepackage{amsmath,amsthm}
% \usepackage{amsfonts} %not required?
% \usepackage{dsfont} %for what?
%unicode-math overwrites the following commands from the mathtools package: \dblcolon, \coloneqq, \Coloneqq, \eqqcolon. Using the other colon-like commands from mathtools will lead to inconsistencies. Plus, Using \overbracket and \underbracke from mathtools package. Use \Uoverbracket and \Uunderbracke for original unicode-math definition.
%use exclusively \mathbf and choose math bold style below.
\usepackage[warnings-off={mathtools-colon, mathtools-overbracket}, bold-style=ISO]{unicode-math}

\defaultfontfeatures{
	Fractions=On,
	Mapping=tex-text% to turn "--" into dashes, useful for bibtex%%
}
\defaultfontfeatures[\rmfamily, \sffamily]{
	Fractions=On,
	Mapping=% to leave " alone (disable the default mapping tex-text; requires loading the font afterwards?)
}
\newfontfamily\xitsfamily{XITS}
\newfontfamily\texgyretermesfamily{TeX Gyre Termes}
\newfontfamily\texgyreherosfamily{TeX Gyre Heros}
\newfontfamily\lmfamily{Latin Modern Roman}
\setmainfont{Latin Modern Roman}
\setsansfont{Latin Modern Sans}
\setmonofont{Latin Modern Mono}
\defaultfontfeatures{}%disable default font features to avoid warnings with math fonts.
\setmathfont{XITS Math}
\setmathfont[range={\mathcal,\mathbfcal},StylisticSet=1]{XITS Math}

\usepackage{cleveref}% cleveref should go "laster" than hyperref
%GRAPHICS
\usepackage{pgf}
\usepackage{pgfplots}
	\usetikzlibrary{matrix,fit,plotmarks,calc,trees,shapes.geometric,positioning,plothandlers}
\pgfplotsset{compat=1.11}
\usepackage{graphicx}

\graphicspath{{graphics/},{graphics-dm/}}
\DeclareGraphicsExtensions{.pdf}

%HACKING
\usepackage{printlen}
\uselengthunit{mm}
% 	\newlength{\templ}% or LenTemp?
% 	\setlength{\templ}{6 pt}
% 	\printlength{\templ}
\usepackage{etoolbox} %for addtocmd command
\usepackage{scrhack}% load at end. Corrects a bug in float package, which is outdated but might be used by other packages
\usepackage{xltxtra} %somebody said that this is loaded by fontspec, but does not seem correct: if not loaded explicitly, does not appear in the log and \showhyphens is not corrected.

%Beamer-specific
%ADD
\usepackage{appendixnumberbeamer}
%\setbeamersize{text margin left=0.1cm, text margin right=0.1cm} 
\setbeamertemplate{navigation symbols}{}
\usetheme{BrusselsBelgium}
\usefonttheme{professionalfonts}


\newcommand{\R}{ℝ}
\newcommand{\N}{ℕ}
\newcommand{\Z}{ℤ}
\newcommand{\card}[1]{\lvert{#1}\rvert}
\newcommand{\powerset}[1]{\mathscr{P}(#1)}%\mathscr rather than \mathcal: scr is rounder than cal (at least in XITS Math).
\newcommand{\suchthat}{\;\ifnum\currentgrouptype=16 \middle\fi|\;}
%\newcommand{\Rplus}{\reels^+\xspace}

\AtBeginDocument{%
	\renewcommand{\epsilon}{\varepsilon}
% we want straight form of \phi for mathematics, as recommended in UTR #25: Unicode support for mathematics.
%	\renewcommand{\phi}{\varphi}
}

% with amssymb, but I don’t want to use amssymb just for that.
% \newcommand{\restr}[2]{{#1}_{\restriction #2}}
%\newcommand{\restr}[2]{{#1\upharpoonright}_{#2}}
\newcommand{\restr}[2]{{#1|}_{#2}}%sometimes typed out incorrectly within \set.
%\newcommand{\restr}[2]{{#1}_{\vert #2}}%\vert errors when used within \Set and is typed out incorrectly within \set.
\DeclareMathOperator*{\argmax}{arg\,max}
\DeclareMathOperator*{\argmin}{arg\,min}


%ARG TH
\newcommand{\AF}{\mathcal{AF}}
\newcommand{\labelling}{\mathcal{L}}
\newcommand{\labin}{\textbf{in}\xspace}
\newcommand{\labout}{\textbf{out}}
\newcommand{\labund}{\textbf{undec}\xspace}
\newcommand{\nonemptyor}[2]{\ifthenelse{\equal{#1}{}}{#2}{#1}}
\newcommand{\gextlab}[2][]{
	\labelling{\mathcal{GE}}_{(#2, \nonemptyor{#1}{\ibeatsr{#2}})}
}
\newcommand{\allargs}{A^*}
\newcommand{\args}{A}
\newcommand{\ar}{a}

%MCDA+Arg
\newcommand{\dm}{d}
\newcommand{\ileadsto}{\leadsto}
\newcommand{\mleadsto}[1][\eta]{\leadsto_{#1}}
\newcommand{\ibeats}{\vartriangleright}
\newcommand{\mbeats}[1][\eta]{\vartriangleright_{#1}}


%MISC
\newcommand{\lequiv}{\Vvdash}
\newcommand{\weightst}{W^{\,t}}

%MCDA classical
\newcommand{\crits}{\mathcal{J}}
\newcommand{\altspace}{\mathbb{A}}
\newcommand{\alts}{A}

%Sorting
\newcommand{\cats}{\mathcal{C}}
\newcommand{\catssubsets}{2^\cats}
\newcommand{\catgg}{\vartriangleright}
\newcommand{\catll}{\vartriangleleft}
\newcommand{\catleq}{\trianglelefteq}
\newcommand{\catgeq}{\trianglerighteq}
\newcommand{\alttoc}[2][x]{(#1 \xrightarrow{} #2)}
\newcommand{\alttocat}[3]{(#2 \xrightarrow{#1} #3)}
\newcommand{\alttoI}{(x \xrightarrow{} \left[\underline{C_x}, \overline{C_x}\right])}
\newcommand{\alttocatdm}[3][t]{\left(#2 \thinspace \raisebox{-3pt}{$\xrightarrow{#1}$}\thinspace #3\right)}
\newcommand{\alttocatatleast}[2]{\left(#1 \thinspace \raisebox{-3pt}{$\xrightarrow[]{≥}$}\thinspace #2\right)}
\newcommand{\alttocatatmost}[2]{\left(#1 \thinspace \raisebox{-3pt}{$\xrightarrow[]{≤}$}\thinspace #2\right)}

\newcommand{\source}{\scriptsize}
\newcommand{\commentOC}[1]{{\selectlanguage{french}{\todo{OC : #1}}}}
%Or: \todo[color=green!40]

%this probably requires outdated float package, see doc KomaScript for an alternative.
% \newfloat{program}{t}{lop}
% \floatname{program}{PM}

%\crefname{axiom}{axiom}{axioms}%might be needed for workaround bug in cref when defining new theorems?

%\ifdefined\theorem\else
%\newtheorem{theorem}{\iflanguage{english}{Theorem}{Théorème}}
%\fi

%which line breaks are chosen: accept worse lines, therefore reducing risk of overfull lines. Default = 200
\tolerance=2000
%accept overfull hbox up to...
\hfuzz=2cm
%reduces verbosity about the bad line breaks
\hbadness 5000
%sloppy sets tolerance to 9999
\apptocmd{\sloppy}{\hbadness 10000\relax}{}{}

% WRITING
%\newcommand{\ie}{i.e.\@\xspace}%to try
%\newcommand{\eg}{e.g.\@\xspace}
%\newcommand{\etal}{et al.\@\xspace}
\newcommand{\ie}{i.e.\ }
\newcommand{\eg}{e.g.\ }
\newcommand{\mkkOK}{\checkmark}%\color{green}{\checkmark}
\newcommand{\mkkREQ}{\ding{53}}%requires pifont?%\color{green}{\checkmark}
\newcommand{\mkkNO}{}%\text{\color{red}{\textsf{X}}}

\makeatletter
\newcommand{\boldor}[2]{%
	\ifnum\strcmp{\f@series}{bx}=\z@
		#1%
	\else
		#2%
	\fi
}
\newcommand{\textstyleElProm}[1]{\boldor{\MakeUppercase{#1}}{\textsc{#1}}}
\makeatother
\newcommand{\electre}{\textstyleElProm{Électre}\xspace}
\newcommand{\electreIv}{\textstyleElProm{Électre Iv}\xspace}
\newcommand{\electreIV}{\textstyleElProm{Électre IV}\xspace}
\newcommand{\electreIII}{\textstyleElProm{Électre III}\xspace}
\newcommand{\electreTRI}{\textstyleElProm{Électre Tri}\xspace}
% \newcommand{\utadis}{\texorpdfstring{\textstyleElProm{utadis}\xspace}{UTADIS}}
% \newcommand{\utadisI}{\texorpdfstring{\textstyleElProm{utadis i}\xspace}{UTADIS I}}

%TODO
% \newcommand{\textstyleElProm}[1]{{\rmfamily\textsc{#1}}} 


\newlength{\GraphsNodeSep}
\setlength{\GraphsNodeSep}{7mm}

% MCDA Drawing Sorting
\newlength{\MCDSCatHeight}
\setlength{\MCDSCatHeight}{6mm}
\newlength{\MCDSAltHeight}
\setlength{\MCDSAltHeight}{4mm}
%separation between two vertical alts
\newlength{\MCDSAltSep}
\setlength{\MCDSAltSep}{2mm}
\newlength{\MCDSCatWidth}
\setlength{\MCDSCatWidth}{3cm}
\newlength{\MCDSEvalRowHeight}
\setlength{\MCDSEvalRowHeight}{6mm}
\newlength{\MCDSAltsToCatsSep}
\setlength{\MCDSAltsToCatsSep}{1.5cm}
\newcounter{MCDSNbAlts}
\newcounter{MCDSNbCats}
\newlength{\MCDSArrowDownOffset}
\setlength{\MCDSArrowDownOffset}{0mm}

\tikzset{/Graphs/dot/.style={
	shape=circle, fill=black, inner sep=0, minimum size=1mm
}}
\tikzset{/MC/D/S/alt/.style={
	shape=rectangle, draw=black, inner sep=0, minimum height=\MCDSAltHeight, minimum width=2.5cm, anchor=north east
}}
\tikzset{MC/D/S/pref/.style={
	shape=ellipse, draw=gray, thick
}}
\tikzset{/MC/D/S/cat/.style={
	shape=rectangle, draw=black, inner sep=0, minimum height=\MCDSCatHeight, minimum width=\MCDSCatWidth, anchor=north west
}}
\tikzset{/MC/D/S/evals matrix/.style={
	matrix, row sep=-\pgflinewidth, column sep=-\pgflinewidth, nodes={shape=rectangle, draw=black, inner sep=0mm, text depth=0.5ex, text height=1em, minimum height=\MCDSEvalRowHeight, minimum width=12mm}, nodes in empty cells, matrix of nodes, inner sep=0mm, outer sep=0mm, row 1/.style={nodes={draw=none, minimum height=0em, text height=, inner ysep=1mm}}
}}

% Beliefs
\tikzset{/Beliefs/D/S/attacker/.style={
	shape=rectangle, draw, minimum size=8mm
}}
\tikzset{/Beliefs/D/S/supporter/.style={
	shape=circle, draw
}}

\newcommand{\tikzmark}[1]{%
	\tikz[overlay, remember picture, baseline=(#1.base)] \node (#1) {};%
}



\setbeamertemplate{headline}[singleline]
\setbeamertemplate{footline}[onlypage]

\title{Syntaxe}
\subject{Java}
\author{Olivier Cailloux}
\institute[LAMSADE]{LAMSADE, Université Paris-Dauphine}
\date{Version du \today}

\begin{document}
\begin{frame}[plain]
	\tikz[remember picture,overlay]{
		\path (current page.south west) node[anchor=south west, inner sep=0] {
			\includegraphics[height=1cm]{LAMSADE95.jpg}
		};
		\path (current page.south) ++ (0, 1mm) node[anchor=south, inner sep=0] {
			\includegraphics[height=9mm]{Dauphine.jpg}
		};
		\path (current page.south east) node[anchor=south east, inner sep=0] {
			\includegraphics[height=9mm]{PSL.png}
		};
	}
	\titlepage
\end{frame}
\addtocounter{framenumber}{-1}

\section{Éléments}
\begin{frame}[fragile]
	\frametitle{Syntaxe élémentaire}
	\begin{itemize}
		\item Instructions terminées par un point virgule
		\item Blocs d’instructions
		\item Commentaires
	\end{itemize}
	\begin{block}{Bloc}
		\begin{lstlisting}
{
	instruction1;
	instruction2;
}
		\end{lstlisting}	
	\end{block}
	\begin{block}{Commentaires}
		\begin{lstlisting}
//commentaire sur une ligne
/* commentaire sur 
        plusieurs lignes */
		\end{lstlisting}
	\end{block}
\end{frame}

\section{Variables}
\begin{frame}[fragile]
	\frametitle{Variables}
	\begin{block}{Déclaration variable}
		\begin{lstlisting}
<type> nom;
		\end{lstlisting}	
	\end{block}
	\begin{block}{Affectation variable}
		\begin{lstlisting}
nom = valeur;
		\end{lstlisting}
	\end{block}
	\begin{block}{Déclaration et affectation variable}
		\begin{lstlisting}
<type> nom = valeur;
		\end{lstlisting}	
	\end{block}
	$⇒$ Essayer avec \texttt{jshell} !
\end{frame}

\begin{frame}
	\frametitle{Types primitifs et valeurs primitives}
	\begin{block}{Types primitifs}
		\begin{itemize}
			\item \texttt{boolean}
			\item \texttt{int} (32 bits)
			\item \texttt{double} (64 bits)
		\end{itemize}
	\end{block}
	\begin{block}{Valeurs primitives}
		\begin{itemize}
			\item \texttt{true}, \texttt{false}
			\item littéral entier : \texttt{156}, \texttt{100\_000}
			\item littéral flottant : \texttt{1.5d}, \texttt{2d} 
		\end{itemize}
	\end{block}
	Une fois qu’une variable a une valeur, on peut l’utiliser
\end{frame}

\begin{frame}[fragile]
	\frametitle{Portée}
	\begin{itemize}
		\item Variables connues en fonction de leur endroit de déclaration
		\item Connues dans les blocs inclus
		\item Inconnues dans les blocs non inclus
	\end{itemize}
	\begin{block}{Exemple}
		\begin{lstlisting}
{
	int a = 3;
	{
		int b = 2;
		// a ? §\onslide<2->{Connu}§
	}
	// a ? §\onslide<3->{Connu}§
	// b ? §\onslide<4->{Inconnu}§
}
		\end{lstlisting}	
	\end{block}
	Bonne pratique : adopter portée la plus étroite possible
\end{frame}

\section{Opérations}
\begin{frame}[fragile]
	\frametitle{Opérateurs unaires}
	Opérateurs unaires :
	\begin{description}
		\item[-] renvoie le négatif
		\item[!] Négation logique
		\item[++] Incrémente un entier (affectation implicite)
		\item[-{}-] Décrémente un entier (affectation implicite)
	\end{description}
	\begin{block}{Exemples}
		\begin{lstlisting}
int x = -3;
int z = -(x - y);
++i;
boolean b = !true;
boolean c = !b;
		\end{lstlisting}	
	\end{block}
\end{frame}

\begin{frame}[fragile]
	\frametitle{Opérations binaires}
	\begin{description}
		\item[+, -, *, /] NB / sur entiers ≠ / sur double : toujours préciser le type pour meilleure lisibilité
		\item[\%] reste de la division entière
		\item[<, >, <=, >=] compare et renvoie un booléen
		\item[==] teste égalité et renvoie un booléen
		\item[!=] teste différence
		\item[\&\&] renvoie vrai ssi les deux opérands sont vrais {\tiny (court-circuite)}
		\item[||] renvoie vrai ssi au moins un opérand est vrai {\tiny (court-circuite)}
		\item[+=, -=] opération puis affectation
	\end{description}
	\begin{block}{Exemple}
		\begin{lstlisting}
boolean c = (!b && (x > y)) || (a != 3)
		\end{lstlisting}	
	\end{block}
\end{frame}

\section{Tableaux}
\begin{frame}[fragile]
	\frametitle{Tableaux}
	\begin{block}{Tableaux}
		\begin{lstlisting}
int[] premiers = {2, 3, 5, 7, 11, 13, 17 };
int x = premiers[0] + 2 * premiers[1];
		\end{lstlisting}	
	\end{block}
	\begin{itemize}
		\item Tableaux à plusieurs dimensions possibles
		\item Dans ce cours on n’utilisera (presque) pas les tableaux
	\end{itemize}
\end{frame}

\section{Branchement}
\begin{frame}[fragile]
	\frametitle{Tests}
	\begin{itemize}
		\item \texttt{if (test) bloc}
		\item \texttt{if (test) bloc else bloc}
		\item \texttt{if (test) bloc else if (test) bloc else bloc}
	\end{itemize}
	\begin{block}{Exemple}
		\begin{lstlisting}
if (value == 3) {
	openGate = true;
} else {
	openGate = false;
}
		\end{lstlisting}	
	\end{block}
	Amélioration ? \pause \texttt{openGate = (value == 3);}
\end{frame}

\begin{frame}[fragile]
	\frametitle{Boucles}
	\begin{block}{While}
		\begin{lstlisting}
while (openGate) {
	§…§
}
do {
	§…§
} while (openGate);
		\end{lstlisting}	
	\end{block}
	\begin{block}{For}
		\begin{lstlisting}
int a = 4;
for (int i=0; i < 3; ++i) {
	a += i;
}
		\end{lstlisting}	
	\end{block}
	Valeur de \texttt{a} ? \pause \texttt{7}
\end{frame}

\section{Méthodes}
\begin{frame}[fragile]
	\frametitle{Méthodes : utilité}
	\begin{itemize}
		\item Tout code doit se trouver dans une méthode
		\item Méthode peut renvoyer une valeur
		\item Peut être réutilisée, clarifie le code
	\end{itemize}
\end{frame}

\begin{frame}[fragile]
	\frametitle{Méthodes : syntaxe}
	\begin{itemize}
		\item \texttt{public static <type de retour> <nom> (<type param1> <nom param1>, …)}
		\item \texttt{return} pour renvoyer valeur (exécution méthode cesse)
		\item \texttt{void} pour indiquer absence de retour
	\end{itemize}
	\begin{block}{Exemple}
		\begin{lstlisting}
public static int weightedSum (int a, int b, 
      int w1, int w2) {
	return a * w1 + b * w2;
}
public static void useWeightedSum () {
	int result = weightedSum(3, 4, 5, 6);
}
		\end{lstlisting}	
	\end{block}
	Valeur de \texttt{result} ? \pause \texttt{39}
\end{frame}

\appendix
\makeatletter
\def\insertframenumber{\@roman\c@framenumber}
\def\inserttotalframenumber{\@roman\c@framenumber}
\makeatother
\AtBeginSection{
}

\clearpage\pdfbookmark{License}{License}
\begin{frame}[plain]
	\frametitle{License}
	This presentation, and the associated \LaTeX{} code, are published under the \href{https://opensource.org/licenses/MIT}{MIT license}. Feel free to reuse (parts of) the presentation, under condition that you cite the author.
	
	Credits are to be given to \href{http://www.lamsade.dauphine.fr/~ocailloux/}{Olivier Cailloux}, Université Paris-Dauphine.
\end{frame}
\addtocounter{framenumber}{-1}
\end{document}

\begin{frame}
	\frametitle{Title}
	\begin{block}{Block}
%		\setlength\abovedisplayskip{1 ex}% reduce space above equations
		\begin{itemize}
			\item Item
		\end{itemize}
	\end{block}
	\begin{itemize}
		\item Item
	\end{itemize}
\end{frame}

